\section{Mid-Term Examination --- Set B}
\label{sec:midterm_set_b}

\ExamHeader{Mid-Term Examination}{B}{60 Minutes}{30 Marks}{Units 1, 2 \& 3 (Languages, Fundamentals \& Hardware)}

\noindent
\textbf{Instructions:} All 30 questions are compulsory. Each question carries 1 mark. Select the single best answer.

\vspace{3mm}

\mcqitem{1}{Which of the following programming languages provides machine independence and human readability?}{Machine Language}{High-Level Language}{Microcode}{Hardware Description}

\mcqitem{2}{A program that translates an entire high-level program into machine code prior to execution is a/an:}{Interpreter}{Compiler}{Assembler}{Text Editor}

\mcqitem{3}{In an assembly instruction \texttt{MOV R1, \#5}, what does \texttt{\#5} represent?}{An opcode}{An immediate operand value}{A memory address pointer}{A label}

\mcqitem{4}{Checking for type compatibility between variables and operators is performed during:}{Lexical Analysis}{Syntax Analysis}{Semantic Analysis}{Code Generation}

\mcqitem{5}{What does the preprocessor in C/C++ handle before compilation?}{Memory allocation}{Header file inclusions and macro replacements}{Linking object files}{Bytecode generation}

\mcqitem{6}{Which tool resolves external library calls (e.g., \texttt{printf}) and merges object files?}{Compiler}{Linker}{Loader}{Debugger}

\mcqitem{7}{What is the primary task of a two-pass assembler during its second pass?}{Building the symbol table}{Generating binary machine instructions}{Checking spelling of comments}{Running the program}

\mcqitem{8}{Testing program behavior with boundary conditions like $hours = 0$ is called:}{Boundary value testing}{Syntax parsing}{Lexical scanning}{Semantic checking}

\mcqitem{9}{Which stage of the CPU instruction cycle determines what operation needs to be performed?}{Fetch}{Decode}{Execute}{Store}

\mcqitem{10}{Who invented the mechanical adding machine known as the Pascaline in 1642?}{Charles Babbage}{Blaise Pascal}{Herman Hollerith}{Alan Turing}

\mcqitem{11}{The Analytical Engine containing concepts of store and mill was designed in the 1830s by:}{Blaise Pascal}{Charles Babbage}{John von Neumann}{Herman Hollerith}

\mcqitem{12}{Herman Hollerith used which technology to automate data compilation for the 1890 US Census?}{Vacuum tubes}{Punched cards}{Transistors}{Integrated circuits}

\mcqitem{13}{First-generation computers (1940--1956) were characterized by:}{Extremely small size and low heat}{Large size, high power consumption, and vacuum tubes}{Silicon microprocessors}{Optical storage}

\mcqitem{14}{Second-generation computers (1956--1963) introduced:}{Integrated circuits}{Transistors and high-level languages like FORTRAN and COBOL}{The internet}{Artificial Intelligence}

\mcqitem{15}{Keyboards, monitors, and multiprogramming operating systems appeared during which generation?}{First generation}{Second generation}{Third generation}{Fifth generation}

\mcqitem{16}{The fourth generation of computers (from 1971) is primarily marked by:}{Vacuum tubes}{Single-chip microprocessors and personal computers}{Relays}{Punched cards}

\mcqitem{17}{Fifth-generation computing focuses on:}{Vacuum tubes}{Artificial Intelligence, parallel processing, and neural networks}{Punched card tabulators}{Analog mechanical gears}

\mcqitem{18}{The computer's ability to perform repetitive operations continuously without fatigue is termed:}{Speed}{Versatility}{Diligence}{Automation}

\mcqitem{19}{Which CPU component performs arithmetic additions and logical comparisons?}{Control Unit}{Arithmetic Logic Unit (ALU)}{Instruction Register}{Address Bus}

\mcqitem{20}{Which system bus transmits physical memory addresses from the CPU to memory chips?}{Data Bus}{Address Bus}{Control Bus}{Serial Bus}

\mcqitem{21}{Which system bus conveys read, write, and interrupt commands between hardware units?}{Data Bus}{Address Bus}{Control Bus}{Power Bus}

\mcqitem{22}{High-performance computers used for weather forecasting, molecular modeling, and physics simulations are:}{Microcontrollers}{Supercomputers}{Workstations}{Embedded systems}

\mcqitem{23}{Large, reliable computers used by banks and government agencies for high-volume transactions are:}{Mainframe computers}{Workstations}{Laptops}{Microcontrollers}

\mcqitem{24}{A dedicated computer built directly into an automobile braking system or microwave oven is a/an:}{Supercomputer}{Embedded computer}{Mainframe}{Server}

\mcqitem{25}{A single-user desktop system with high-end CPU and graphics for CAD and 3D engineering is a:}{Mainframe}{Workstation}{Embedded system}{Supercomputer}

\mcqitem{26}{Which of the following is non-volatile primary memory that stores startup firmware (BIOS/UEFI)?}{RAM}{ROM}{L1 Cache}{L2 Cache}

\mcqitem{27}{Dynamic RAM (DRAM) is used as main system memory because:}{It uses flip-flops and needs no refresh}{It provides high density at moderate cost, though it requires periodic electrical refreshing}{It is non-volatile}{It is made of magnetic tape}

\mcqitem{28}{Static RAM (SRAM) is used in CPU cache memories primarily because:}{It is very inexpensive}{It is faster than DRAM and does not require periodic refreshing}{It has higher capacity than hard drives}{It retains data when unpowered}

\mcqitem{29}{A mechanical Hard Disk Drive (HDD) is more prone to physical shock damage than an SSD because:}{It uses high-voltage lasers}{It contains moving mechanical read/write heads and rotating platters}{It is made of silicon chips}{It has no moving parts}

\mcqitem{30}{An input device that converts physical sound waves into digital electrical signals is a:}{Speaker}{Microphone}{Monitor}{Printer}

\newpage
\subsection*{Answer Key \& Explanations --- Mid-Term (Set B)}

\begin{answerkeytable}{Mid-Term Set B --- Answer Key}
\begin{tabular}{*{10}{c}}
\toprule
\textbf{Q} & \textbf{Ans} & \textbf{Q} & \textbf{Ans} & \textbf{Q} & \textbf{Ans} & \textbf{Q} & \textbf{Ans} & \textbf{Q} & \textbf{Ans} \\
\midrule
1 & \textbf{B} & 7 & \textbf{B} & 13 & \textbf{B} & 19 & \textbf{B} & 25 & \textbf{B} \\
2 & \textbf{B} & 8 & \textbf{A} & 14 & \textbf{B} & 20 & \textbf{B} & 26 & \textbf{B} \\
3 & \textbf{B} & 9 & \textbf{B} & 15 & \textbf{C} & 21 & \textbf{C} & 27 & \textbf{B} \\
4 & \textbf{C} & 10 & \textbf{B} & 16 & \textbf{B} & 22 & \textbf{B} & 28 & \textbf{B} \\
5 & \textbf{B} & 11 & \textbf{B} & 17 & \textbf{B} & 23 & \textbf{A} & 29 & \textbf{B} \\
6 & \textbf{B} & 12 & \textbf{B} & 18 & \textbf{C} & 24 & \textbf{B} & 30 & \textbf{B} \\
\bottomrule
\end{tabular}
\end{answerkeytable}

\begin{expbox}[Technical Solutions --- Mid-Term Set B]
\begin{itemize}[leftmargin=5mm]
    \item \textbf{Q.4 (C):} Semantic analysis checks for type mismatches, variable declarations, and scope rules.
    \item \textbf{Q.9 (B):} The Decode stage determines the operation and required operands from the fetched opcode.
    \item \textbf{Q.20 (B):} The Address Bus carries the physical memory address from the CPU to memory chips.
    \item \textbf{Q.27 (B):} DRAM stores bits in tiny capacitors that lose charge, requiring dynamic refreshing.
    \item \textbf{Q.28 (B):} SRAM uses transistor flip-flops, making it much faster than DRAM and ideal for CPU cache.
\end{itemize}
\end{expbox}

\newpage
